\documentclass[10pt,twoside]{article}
\usepackage[utf8]{inputenc}
\usepackage{amsmath}
\usepackage{amsfonts}
\usepackage{amssymb}
\usepackage[spanish,es-noshorthands]{babel}
\usepackage[T1]{fontenc}
\usepackage{lmodern}
\usepackage{graphicx,hyperref}
\usepackage{tikz,pgf}
\usepackage{multicol}
\usepackage{subfig}
\usepackage[papersize={6.5in,8.5in},width=5.5in,height=7in]{geometry}
\usepackage{fancyhdr}
\pagestyle{fancy}
\fancyhead[LE]{\includegraphics[height=12pt]{Images/logo-colegio.png} Cálculo $11^{\circ}$}
\fancyhead[RE]{}
\fancyhead[RO]{\textit{Germ\'an Avenda\~no Ram\'irez, Lic. U.D., M.Sc. U.N.}}
\fancyhead[LO]{}

\author{Germ\'an Avenda\~no Ram\'irez, Lic. U.D., M.Sc. U.N.}
\title{\begin{minipage}{.2\textwidth}
\includegraphics[height=1.75cm]{Images/logo-colegio.png}\end{minipage}
\begin{minipage}{.55\textwidth}
\begin{center}
Animaplano 1, $10^{\circ}$  \\
Cálculo $11^{\circ}$
\end{center}
\end{minipage}\hfill
\begin{minipage}{.2\textwidth}
\includegraphics[height=1.75cm]{Images/logo-sed.png} 
\end{minipage}}
\date{}
\begin{document}
\maketitle
Nombre: \hrulefill Curso: \underline{\hspace*{44pt}} Fecha: \underline{\hspace*{2.5cm}}
\section*{Cuestionario}
Responda haciendo los procedimientos al frente de cada pregunta
\begin{enumerate}
 \item $8^{2}=4^{3}=2^{6}=$
 \item Si $2^{x}=64$, entonces $x+61=$?
 \item $10^{2}-6^{2}-4^{2}+3^{2}=$
 \item La media de los datos \{63, 64, 65, 66, 67\} es:
 \item $22xy-3xy+35xy=\underline{\hspace{24pt}}xy$
 \item Calcule $\log_{2}256^{8}=$
 \item El 25\% de 292 es:
 \item $9^{4}/9^{2}-10^{0}$ menos raíz cuadrada de 81 =
 \item (raíz cuadrada de 64 $\times$ raíz cuadrada de 64) menos 4
 \item la pendiente de la recta $y=45+81x$ es?
 \item El vigésimo término de la progresión aritmética 26, 29, 32, 35, \dots es:
 \item $10^{2}-10^{2}/4=$
 \item Si 3/2=n/14, entonces $(n\times4)+9=$
 \item Si $sen(n-17)=\frac{\sqrt{3}}{2}$, entonces $n=$?
 \item Si 15 es el primer término de una progresión aritmética cuya diferencia entre dos términos consecutivos es 8, entonces el 10 término es:
 \item La imagen de 4 en la función $f(x)=x^{3}+2x+6$ es:
 \item El tercer término de una progresión geométrica cuya razón es 2 y cuyo primer término es 17 es:
 \item $5!-6^{2}-5$
 \item Si $2x=10$, entonces $x^{2}+3x+18=$?
 \item $(14-18+2)^{2})\cdot(4+12-4)$
 \item $\log_{8}8+\log_{2}64^{6}=$
 \item $4^{2x}=2^{40}$, entonces $x+35=$?
 \item $3^{3x}=27$, entonces $x+44=$?
 \item Si $4x=24$, entonces $x^{2}=$?
 \item Halle $4!+2=$?
 \item El doceavo término de una progresión aritmética cuyos primeros términos son -14, --11, -8, -5, \ldots es:
 \item $(\sqrt[3]{64})^{2}=$
 \item El cuarto término de una progresión geométrica cuyo primer término es $\frac{1}{3}$ y cuya razón es 3.
 \item El tercer número primo
 \item Número de dos cifras cuyo producto es 4 y cuya suma es 5.
 \item Si el perímetro del rectángulo \tikz \draw (0,0)--node{$x$}(3,0)--(3,1)--(0,1)--cycle; mide 40, entonces $x+10=$?
\end{enumerate}

\end{document}
